Esta guía le muestra, paso a paso y con imágenes, cómo instalar \textbf{ConcretoLab Desktop} en una computadora con
Windows: desde la descarga del instalador hasta el primer inicio del programa. \textbf{No necesita instalar Java ni
ningún otro programa}: el instalador ya trae todo lo necesario. Solo se necesita internet para descargarlo, y no se
requieren permisos de administrador.

\begin{nota}[La instalación en 6 pasos]
\begin{enumerate}
	\item Descargue \ruta{\archivoinstalador} desde el enlace oficial (sección~\ref{sec:descarga}).
	\item Haga doble clic en el archivo descargado.
	\item Si Windows muestra el aviso «Windows protegió su PC», pulse \boton{Más información} y luego
	      \boton{Ejecutar de todas formas} (página \pageref{sec:smartscreen}).
	\item En el asistente pulse \boton{Siguiente} hasta llegar a \boton{Instalar}.
	\item Cuando termine, pulse \boton{Finalizar}.
	\item Abra ConcretoLab desde el menú Inicio o desde el acceso del escritorio.
\end{enumerate}
\end{nota}

\textbf{Cómo leer las imágenes.} En cada captura, un número en un círculo rojo (por ejemplo \marca{1}) y un recuadro rojo
señalan lo que hay que pulsar, en ese orden. Los botones de las pantallas se escriben así: \boton{Siguiente}. Las teclas
del teclado, así: \tecla{Enter}.

\textbf{A quién va dirigida.} A cualquier persona que vaya a usar ConcretoLab: no se necesitan conocimientos de
informática. La última sección resuelve los problemas más frecuentes y el anexo explica, para el personal de soporte
técnico, cómo comprobar el archivo y cómo instalar sin ventanas.
